\chapter{Ejemplos}
\label{cap:ejemplos}

\section{Dibujos y colores}
\label{sec:dibujos}

\begin{figure}[ht!]
\centering
\begin{tikzpicture}
\node (0) [vertex,label=180:$v_1$] at (0,0){};
\node (1) [vertex,label=90:$v_2$]  at (1,0){};
\begin{scope}[xshift=2cm]
\node (2) [vertex,label=90:$v_3$]  at (270:1){};
\node (3) [vertex,label=0:$v_4$]   at (0:1){};
\node (4) [vertex,label=90:$v_5$]  at (90:1){};
\end{scope}

\draw [myloop,in=60,out=120,looseness=12]
      (0) to node[above]{$e_1$} (0);
\draw [myloop,in=240,out=300,looseness=12]
      (0) to node[below]{$e_2$} (0);
\draw [myloop,in=240,out=300,looseness=12]
      (2) to node[below]{$e_9$} (2);

\draw [edge]               (0) to node [above] {$e_3$} (1);
\draw [edge]               (1) to node [below] {$e_4$} (2);
\draw [edge]               (2) to node [below] {$e_5$} (3);
\draw [edge,bend right=20] (3) to node [above] {$e_6$} (4);
\draw [edge,bend left=20]  (3) to node [below] {$e_7$} (4);
\draw [edge]               (4) to node [above] {$e_8$} (1);


% Componente derecha
\begin{scope}[xshift=6cm]
\node (5) [vertex,label=90:$v_6$]    at (0:1.5){};
\node (6) [vertex,label=120:$v_7$]   at (120:1.5){};
\node (7) [vertex,label=240:$v_8$]   at (240:1.5){};
\node (8) [vertex,label=0:$v_9$]     at (3,1.2){};
\node (9) [vertex,label=0:$v_{10}$]  at (3,-1.2){};

\draw [edge,bend right=10] (5) to node [above] {$e_{10}$} (6);
\draw [edge,bend left=10]  (5) to node [below] {$e_{11}$} (6);
\draw [edge,bend right=10] (6) to node [left]  {$e_{12}$} (7);
\draw [edge,bend left=10]  (6) to node [right] {$e_{13}$} (7);
\draw [edge,bend right=10] (7) to node [below] {$e_{14}$} (5);
\draw [edge,bend left=10]  (7) to node [above] {$e_{15}$} (5);
\draw [edge]               (5) to node [above] {$e_{16}$} (8);
\draw [edge]               (5) to node [below] {$e_{17}$} (9);
\end{scope}

\end{tikzpicture}
\caption{El diagrama de una gr\'afica con lazos y
aristas m\'ultiples.}
\label{fig:grafica}
\end{figure}

\begin{figure}[ht!]
\begin{tikzpicture}
%%%%%%%%%%%%%%%%%%%%%%%%%%%%%%%%%%%%%%%%%%%%%%%%
%%%%%%%%%%         Empty Graph        %%%%%%%%%%
%%%%%%%%%%%%%%%%%%%%%%%%%%%%%%%%%%%%%%%%%%%%%%%%
\begin{scope}
% if label is needed -> label={(360/4)*\i}:$\i$
\foreach \i in {0,...,3}
	\node (\i) [vertex,fill=cyan] at ({(360/4)*\i}:1){};

\node (L) at (-1,1){$G_1$};
\end{scope}

%%%%%%%%%%%%%%%%%%%%%%%%%%%%%%%%%%%%%%%%%%%%%%%%
%%%%%%%%%%       Complete Graph       %%%%%%%%%%
%%%%%%%%%%%%%%%%%%%%%%%%%%%%%%%%%%%%%%%%%%%%%%%%
\begin{scope}[xshift=3.1cm,yshift=-0.3cm]
\node (4) [vertex,fill=verdeAzulado] at (0,0){};
\foreach \i in {0,1,2}
	\node (\i) [vertex,fill=verdeAzulado] at ({90+(360/3)*\i}:1.2){};

\foreach \i in {0,1,2}
	\draw [edge] let \n1={int(mod(\i+1,3))} in (\i) to (\n1);
\foreach \i in {0,1,2}
	\draw [edge] (\i) to (4);

\node (L) at (-1,1){$G_2$};
\end{scope}


%%%%%%%%%%%%%%%%%%%%%%%%%%%%%%%%%%%%%%%%%%%%%%%%
%%%%%%%%%%       Bipartite Graph      %%%%%%%%%%
%%%%%%%%%%%%%%%%%%%%%%%%%%%%%%%%%%%%%%%%%%%%%%%%
\begin{scope}[xshift=6.2cm,yshift=0cm]
\foreach \i in {0,2,4}
	\node (\i) [vertex,fill=red]  at ({(360/6)*\i}:1){};
\foreach \i in {1,3,5}
	\node (\i) [vertex,fill=blue] at ({(360/6)*\i}:1){};

\foreach \i in {0,...,5}
	\draw [edge] let \n1={int(mod(\i+1,6))} in (\i) to (\n1);

\node (L) at (-1,1){$G_3$};
\end{scope}


%%%%%%%%%%%%%%%%%%%%%%%%%%%%%%%%%%%%%%%%%%%%%%%%
%%%%%%%%%%            Star            %%%%%%%%%%
%%%%%%%%%%%%%%%%%%%%%%%%%%%%%%%%%%%%%%%%%%%%%%%%
\begin{scope}[xshift=9.3cm]
\node (6) [vertex,fill=violet] at (0,0){};
\foreach \i in {0,...,4}
	\node (\i) [vertex,fill=green] at ({90+(360/5)*\i}:1){};

\foreach \i in {0,...,4}
	\draw [edge] (\i) to (6);

\node (L) at (-1,1){$G_4$};
\end{scope}


%%%%%%%%%%%%%%%%%%%%%%%%%%%%%%%%%%%%%%%%%%%%%%%%
%%%%%%%%%%           Split            %%%%%%%%%%
%%%%%%%%%%%%%%%%%%%%%%%%%%%%%%%%%%%%%%%%%%%%%%%%
\begin{scope}[xshift=12.4cm]
\foreach \i in {0,2,4}
	\node (\i) [vertex,fill=black] at ({30+(360/6)*\i}:0.7){};
\foreach \i in {1,3,5}
	\node (\i) [vertex] at ({30+(360/6)*\i}:1.4){};

\foreach \i in {0,...,5}
	\draw [edge] let \n1={int(mod(\i+1,6))} in (\i) to (\n1);
\foreach \i in {0,2,4}
	\draw [edge] let \n1={int(mod(\i+2,6))} in (\i) to (\n1);

\node (L) at (-1,1){$G_5$};
\end{scope}

\end{tikzpicture}
\caption{Ejemplos de gr\'aficas vac\'ia, completa,
bipartita, bipartita completa y escindible.}
\label{fig:fam1}
\end{figure}
